\section{Weak Formulations}

Finally, we turn to the {\em weak formulation} of Signorini's problem. Unlike the weak formulation of conventional simple boundary value problems, here we do not end up with a variational equation, but with a variational inequality. This is derived by a first-order optimality condition on the functional $\cJ$ under the constraints of $K$, i.e. a condition on the {\em Fréchet derivative} of $\cJ$. To fix notation, we start with a definition.

\begin{definition}[Fréchet derivative]\label{def:frechetdiff}
	Let $X$ and $Y$ be normed vector spaces and let
	\begin{equation*}
		f:U\subseteq X\to Y
	\end{equation*}
	be a function defined on an open subset $U$ of $X$. Then $f$ is called {\tt Fréchet differentiable} at $x \in U$ if there exists a bounded linear operator $A(x):X\to Y$ such that
	\begin{equation*}
		\lim_{\|h\|_X\to 0} \frac{\|f(x+h)-f(x)-A(x)[h]\|_Y}{\|h\|_X}=0.
	\end{equation*}
	If $A(x)$ exists, it is unique and is called the {\tt Fréchet derivative} of $f$ at $x$. We write
	\begin{equation*}
		f^\prime(x)[h]=A(x)[h]\qquad \forall\, h\in X.
	\end{equation*}
\end{definition}

Finally, we prove Fréchet differentiability of the energy functional
\begin{equation*}
	\cJ(v) := \frac{1}{2}\,a(v,v)-\ell(v),\qquad v\in V
\end{equation*}
and provide an explicit formula for its Fréchet derivative.

\begin{lemma}[Fréchet differentiability of $\cJ$]\label{lem:Jdiff}
	The functional $\cJ$ defined in \eqref{eq:energyfunc} is Fréchet differentiable on $V$. Its derivative at $u\in V$ is given by
	\begin{equation*}
		\cJ'(u)[h] = a(u,h)-\ell(h)\qquad \forall\, h\in V.
	\end{equation*}
\end{lemma}

\begin{proof}
	Let $u,k\in V$. Using bilinearity of $a$ and linearity of $\ell$, we obtain
	\begin{align*}
		\cJ(u+k)-\cJ(u)
		 & =\frac12\bigl(a(u+k,u+k)-a(u,u)\bigr)-\ell(k)            \\
		 & =\frac12\bigl(a(u,u)+2a(u,k)+a(k,k)-a(u,u)\bigr)-\ell(k) \\
		 & = a(u,k)-\ell(k)+\frac12\,a(k,k).
	\end{align*}
	Hence
	\begin{equation*}
		\cJ(u+k)-\cJ(u)-\bigl(a(u,k)-\ell(k)\bigr)
		=\frac12\,a(k,k).
	\end{equation*}
	Define $A(u):V\to\mathbb R$ by
	\begin{equation*}
		A(u)[h]:=a(u,h)-\ell(h).
	\end{equation*}
	Since $a$ and $\ell$ are bounded, $A(u)$ is linear and bounded. Moreover,
	\begin{equation*}
		\frac{\bigl|\cJ(u+k)-\cJ(u)-A(u)[k]\bigr|}{\|k\|_{1,\Omega}}
		=\frac{1}{2}\,\frac{|a(k,k)|}{\|k\|_{1,\Omega}}
		\le \frac12\,C_{\mathcal A}\,\|k\|_{1,\Omega}\xrightarrow{\ \|k\|_{1,\Omega}\to 0\ }0,
	\end{equation*}
	by boundedness of $a$, (see Proposition \ref{pr:aproperties}). Therefore, $\cJ$ is Fréchet differentiable at $u$ with
	\begin{equation*}
		\cJ^\prime(u)[h]=A(u)[h]=a(u,h)-\ell(h)
		\qquad \forall\, h\in V.
	\end{equation*}
\end{proof}

\subsection{Variational Inequality as an Optimality Condition}

\begin{proposition}[First-order optimality condition on convex sets]\label{pr:FOC}
	A function $u\in K$ minimizes $\cJ$ over $K$ if and only if it satisfies the variational inequality
	\begin{equation}\label{eq:firstorderopt}
		\cJ^\prime(u)[v-u]\ge 0\qquad\forall\, v\in K.
	\end{equation}
\end{proposition}
\begin{proof}
	Assume first that $u\in K$ minimizes $\cJ$ over $K$ and let $v\in K$.
	By convexity of $K$ we have $u+t(v-u)\in K$ for all $t\in[0,1]$.
	Define $\phi(t):=\cJ(u+t(v-u))$. Then $\phi$ attains its minimum at $t=0$, hence $\phi'(0^+)\ge 0$, where $\phi'(0^+)$ denotes the right-sided derivative at $t=0$. Using the chain rule and Fréchet differentiability of $\cJ$ we obtain
	\begin{equation*}
		\phi'(0^+)=\cJ'(u)[v-u]\ge 0.
	\end{equation*}
	Conversely, assume \eqref{eq:firstorderopt}. Let $v\in K$ be arbitrary and set $h:=v-u$. By Taylor expansion of the quadratic functional $\cJ$ we have
	\begin{equation*}
		\cJ(u+h)-\cJ(u)=\cJ'(u)[h]+\frac12\,a(h,h).
	\end{equation*}
	By \eqref{eq:firstorderopt}, $\cJ'(u)[h]\ge 0$, and since $a(h,h)\ge 0$ we conclude $\cJ(v)\ge \cJ(u)$ for all $v\in K$. Hence $u$ is a minimizer of $\cJ$ over $K$.
\end{proof}

Combining Lemma \ref{lem:Jdiff} and Proposition \ref{pr:FOC}, $u\in K$ solves Problem \ref{pr:signorini1} if and only if
\begin{equation*}
	\cJ^\prime(u)[v-u] = a(u,v-u)-\ell(v-u)\ge 0 \qquad \forall\, v\in K.
\end{equation*}
This is the first weak formulation of Signorini's problem as a variational inequality.

\begin{problem}[Signorini's Problem -- Weak Formulation 1]\label{pr:signorini2}
\begin{equation*}
	\text{\em Find } u\in K\text{\em\enspace such that:}\qquad a(u,v-u)\ge \ell(v-u)\qquad \forall\, v\in K.
\end{equation*}
\end{problem}

Well-posedness of Problem \ref{pr:signorini2} is already ensured, since it is equivalent to the minimization Problem \ref{pr:signorini1} by Proposition \ref{pr:FOC} in the sense that each solution of one of the problems is also a solution of the other. The solution is unique due to Corollary \ref{cor:wellposed} with the stability estimate \eqref{eq:ustab}. Another equivalent formulation is the following.

\begin{problem}[Signorini's Problem -- Weak Formulation 2]\label{pr:signorini3}
\begin{equation*}
	\text{\em Find } u\in K\text{\em\enspace such that:}\qquad\begin{cases}
		a(u,u) = \ell(u), &                   \\[2pt]
		a(u,v)\ge \ell(v) & \forall\, v\in K.
	\end{cases}
\end{equation*}
\end{problem}

\begin{proposition}[Equivalence of Problem \ref{pr:signorini2} and Problem \ref{pr:signorini3}]\label{pr:equivalence23}
	Any solution $u\in K$ to Problem \ref{pr:signorini2} is also a solution to Problem \ref{pr:signorini3} and vice versa.
\end{proposition}
\begin{proof}
	Let $u\in K$ solve Problem \ref{pr:signorini3}. Then
	\begin{equation*}
		a(u,v-u) = a(u,v)-a(u,u)  = a(u,v)-\ell(u)\ge \ell(v)-\ell(u) = \ell(v-u)\qquad\forall\,v\in K.
	\end{equation*}
	Assuming that $u\in K$ solves Problem \ref{pr:signorini2} and testing $v=0\in K$ yields
	\begin{equation*}
		a(u,-u)\ge \ell(-u)\qquad\Longleftrightarrow\qquad -a(u,u)\ge -\ell(u),
	\end{equation*}
	which directly implies
	\begin{equation*}
		a(u,u)\le \ell(u).
	\end{equation*}
	However, testing with $v=2u\in K$ yields $a(u,u)\ge \ell(u)$ and thus $a(u,u) = \ell(u)$. Furthermore, Problem \ref{pr:signorini2} reads
	\begin{equation*}
		a(u,v-u) = a(u,v)-a(u,u)\ge \ell(v-u) = \ell(v)-\ell(u)\qquad\forall\,v\in K,
	\end{equation*}
	which directly implies
	\begin{equation*}
		a(u,v)\ge a(u,u)+\ell(v)-\ell(u)\qquad\forall\,v\in K.
	\end{equation*}
	Plugging in $a(u,u) = \ell(u)$ yields $a(u,v)\ge \ell(v)$ for all $v\in K$.
\end{proof}

\subsection{Lipschitz Dependence of the Displacement on the Forces}

We close this section by providing another stability result on the dependence of solutions to varying source terms as input data. This result is very similar to the stability estimate provided in Corollary \ref{cor:wellposed}. Its proof, however, is much more convenient at this point, where the weak formulations are available.

\begin{proposition}[Lipschitz dependence on the data]\label{pr:lipschitz}
	For $i\in\{1,2\}$, let $f_i\in L^2(\Omega;\mathbb R^d)$ and $F_i\in L^2(\Gamma_N;\mathbb R^d)$, and let $\ell_i\in V'$ be the corresponding load functional, i.e.
	\begin{equation*}
		\ell_i(v):=\int_\Omega f_i\cdot v\;\dx+\int_{\Gamma_N} F_i\cdot \gamma v\;\ds,
		\qquad v\in V.
	\end{equation*}
	Moreover, let $u_i\in K$ be the unique solution of Problem \ref{pr:signorini2} associated with $\ell_i$, i.e.
	\begin{equation*}
		a(u_i,v-u_i)\ge \ell_i(v-u_i)\qquad \forall\, v\in K.
	\end{equation*}
	Then the solution depends Lipschitz continuously on the load functional, i.e.
	\begin{equation}\label{eq:lipschitzabstract}
		\|u_1-u_2\|_{1,\Omega}\le \frac{C_K^2}{\alpha_\mathcal{A}}\,\|\ell_1-\ell_2\|_{V'}.
	\end{equation}
	In particular, it depends Lipschitz continuously on the data $(f_1, F_1)$ and $(f_2, F_2)$ in the sense that
	\begin{equation}\label{eq:lipschitzloads}
		\|u_1-u_2\|_{1,\Omega}
		\le \frac{C_K^2\,\max\{1,C_\mathrm{tr}\}}{\alpha_\mathcal{A}}
		\Bigl(\|f_1-f_2\|_{0,\Omega}+\|F_1-F_2\|_{0,\Gamma_N}\Bigr).
	\end{equation}
\end{proposition}

\begin{proof}
	For $u_1$, choose $v=u_2\in K$ in Problem \ref{pr:signorini2}. Then $a(u_1,u_2-u_1)\ge \ell_1(u_2-u_1)$. Similarly, choosing $v=u_1\in K$ for $u_2$ yields $a(u_2,u_1-u_2)\ge \ell_2(u_1-u_2)$. Adding both inequalities and using symmetry of $a$, we obtain
	\begin{equation*}
		a(u_1-u_2,u_1-u_2)\le (\ell_1-\ell_2)(u_1-u_2).
	\end{equation*}
	By Proposition \ref{pr:aproperties}, the bilinear form $a$ is $V$-elliptic with constant $\alpha_\cA/C_K^2$. Hence
	\begin{equation*}
		\frac{\alpha_\mathcal{A}}{C_K^2}\,\|u_1-u_2\|_{1,\Omega}^2
		\le a(u_1-u_2,u_1-u_2)
		\le (\ell_1-\ell_2)(u_1-u_2)
		\le \|\ell_1-\ell_2\|_{V'}\,\|u_1-u_2\|_{1,\Omega}.
	\end{equation*}
	If $u_1=u_2$, there is nothing to prove. Otherwise, dividing by $\|u_1-u_2\|_{1,\Omega}$ yields \eqref{eq:lipschitzabstract}. Moreover,
	\begin{equation*}
		(\ell_1-\ell_2)(v)
		=\int_\Omega (f_1-f_2)\cdot v\;\dx+\int_{\Gamma_N} (F_1-F_2)\cdot \gamma v\;\ds,
		\qquad v\in V.
	\end{equation*}
	Therefore, Proposition \ref{pr:aproperties} applied to the functional $\ell_1-\ell_2$ implies
	\begin{equation*}
		\|\ell_1-\ell_2\|_{V'}
		\le \max\{1,C_\mathrm{tr}\}
		\Bigl(\|f_1-f_2\|_{0,\Omega}+\|F_1-F_2\|_{0,\Gamma_N}\Bigr).
	\end{equation*}
	Combining this with \eqref{eq:lipschitzabstract}, we obtain \eqref{eq:lipschitzloads}.
\end{proof}